\section{静态树上问题}

\section{链剖分}

%TODO

链剖分把树的\textbf{点集}划分为若干条直上直下的链，两个端点在同一条链中的边称为链内边/重边/实边\graffiti{注意这些边并没有一个公认的命名。}，不在同一条链中的边称为链间边/轻边/虚边。也可以用另一种方式理解这个链划分：每个结点选择至多一个儿子作为特殊儿子，对应的边为链内边，其它边为链间边。根据选择儿子的方式不同，可以得到以下几种链剖分：

\begin{itemize}
    \item 重链剖分：每个结点选择子树大小最大的儿子作为特殊儿子，称为重儿子。此时，链内边称为重边，其它边称为轻边，划分出的链称为重链。\graffiti{由于绝大多数情况下使用到的链剖分都是重链剖分，人们常常也会把重链剖分的术语用到其他链剖分上。}
    \item 长链剖分：每个结点选择子树深度最大（指从根到最深叶子的长度）的儿子作为特殊儿子，称为长儿子/重儿子。此时，链内边称为长边/重边，其它边称为短边/轻边，划分出的链称为长链。
    \item LCT 中使用的动态链剖分：每个结点选择任意一个儿子作为特殊儿子，称为实儿子。此时，链内边称为实边，其它边称为虚边，划分出的链称为实链。
\end{itemize}

在下文中，我们称链剖分中的一条链为\textbf{实链}，而链则用于指代普通的树上路径。$\mathrm{son}_u$ 表示结点 $u$ 的特殊儿子，$\mathrm{child}_u$ 表示结点 $u$ 的所有儿子组成的集合。一条实链中最浅的结点称为这条实链的链顶，$u$ 所在实链的链顶记作 $\mathrm{top}[u]$。

链剖分的基本性质是每条链间边的较深结点都是一条实链的链顶；反过来，一条实链的链顶要么是根，要么是某条链间边的较深结点。不过，需要注意链间边的较浅结点未必是一条实链的链底。根据这个性质，有以下结果成立：
$$
\sum_{u}\sum_{v\in \mathrm{child}_u\backslash \{\mathrm{son}_u\}} v\text{ 到所在实链底的距离}\le \sum_{\text{链 }C}|C|=n.
$$
这是分析复杂度时常用的一个结论。
\begin{exercise}
    令 $l_u$ 为 $u$ 到其子树内最深叶子的长度，即长链剖分的划分依据。在长链剖分中，证明
    $$
    \sum_{u}\sum_{v\in \mathrm{child}_u\backslash \{\mathrm{son}_u\}} l_v\le n.
    $$
    这是长链剖分优化 DP 的一个依据。
\end{exercise}

链剖分的一个简单应用是寻找同一条实链上某个深度的点。
\begin{example}
    给定一个链剖分，$O(1)$ 回答从 $u$ 出发沿着其所在链向上/下走 $k$ 步到达的点，或者报告离开了所在链。
\end{example}
\begin{exercise}[$k$ 级祖先]
    给一棵有根树，多次询问一个点的 $k$ 级祖先。离线 $O(n)-O(1)$，在线 $O(n\log n)-O(1)$。
\end{exercise}
\begin{exercise}
    给一棵有根树，儿子有序排列。多次询问从某个点出发，每次都走最小的儿子，走 $k$ 步后到达的点。保证询问点 $u$ 子树内任何叶子到 $u$ 的距离都至少为 $k$
\end{exercise}


任何一条直上直下的链都可以自下而上地分解为如下结构：（实链前缀、链间边）重复若干次，再加上一个实链的一段。如果是到根的链，则最后一段也是一个实链前缀。我们可以对每条实链维护区间数据结构，从而支持一些链信息查询。我们也可以利用实链快速上跳来计算 LCA。

\begin{descbox}{算法}[链剖分计算 LCA]
    在 $u,v$ 处于同一条实链前，反复执行以下过程：每次从 $u,v$ 中选择链顶深度较大的一个，不妨设为 $u$，将 $u$ 跳到 $fa[\mathrm{top}[u]]$。对于同一条实链上的 $u,v$，LCA 即为深度较小的那个。复杂度为 $u,v$ 路径上链间边的个数。
\end{descbox}

\begin{exercise}
    对于线段树可以直接 $O(\log n)$ 维护的静态区间信息查询，设计一个 $O(u,v\ \text{路径上的链间边个数}+\log n)$ 的算法来支持树路径上的相同信息查询。
\end{exercise}
\begin{exercise}
    对于线段树可以直接 $O(\log n)$ 维护的带修区间信息查询，设计一个 $O(u,v\ \text{路径上的链间边个数}\log n)$ 的算法来支持树路径上的相同修改查询。
\end{exercise}
\begin{remark}
    链剖分+线段树并不是一个完全平衡的结构，这导致复杂度有一个额外的 $\log$ 因子。可以使用全局平衡二叉树来做到 $O(\log n)$ 修改查询。
\end{remark}

显然，链剖分的复杂度分析依赖于链间边的个数。不同的链剖分方式会导致不同的复杂度。

\begin{lemma}
    在重链剖分中，任意根到点的路径上经过的轻边数不超过 $\log_2 n$。
\end{lemma}
\begin{proof}
    从上往下每经过一条轻边，子树大小至少翻倍。
\end{proof}
\begin{exercise}
    证明长链剖分中，任意根到点的路径上经过的短边数为 $O(\sqrt{n})$。
\end{exercise}

可以配合 DFS 序来同时支持维护子树信息。在 DFS 时，我们优先递归 $\mathrm{son}_u$，这样就能保证一条重链在 DFS 序中是连续的。按照这种 DFS 序构建的线段树能够同时支持重链和子树的修改查询。我们也可以设计其他的结点排序方式来支持更复杂的询问。

\begin{example}[毛毛虫剖分]
    对于路径 $P_{uv}$，$\{x:d(x,P_{uv})\le 1\}$ 称为 $u,v$ 的毛毛虫。构造一个结点排序方式，使得子树可以划分为 $O(1)$ 个连续区间，且一个毛毛虫可以划分为 $O(\log n)$ 个连续区间。
\end{example}
\begin{solution}
    重链剖分，然后按照如下顺序构造顺序：写下所有儿子，递归重儿子，然后递归所有轻儿子。这样，对于每条重链 $C$，设其自上而下为 $u_1,\cdots,u_k$，则递归到 $u_1$ 时会连续写下 $\mathrm{child}_{u_1},\mathrm{child}_{u_2},\cdots$。也就是说一条重链对应的毛毛虫中，除了链顶之外都是连续的。同理，对于子树，除了子树根以外都是连续的。
\end{solution}
\begin{exercise}[$k$-毛毛虫剖分]
    对于路径 $P_{uv}$，$\{x:d(x,P_{uv})\le k\}$ 称为 $u,v$ 的 $k$-毛毛虫。构造一个结点排序方式，使得子树可以划分为 $O(k)$ 个连续区间，且一个 $k$-毛毛虫可以划分为 $O(k\log n)$ 个连续区间。
\end{exercise}

\section{Top Tree}

% TODO